\documentclass[11pt]{article}
\usepackage[T1]{fontenc}
\usepackage{lmodern,microtype,amsmath,amssymb,amsthm,booktabs}
\usepackage[margin=1in]{geometry}
\usepackage[hidelinks]{hyperref}
\newtheorem{proposition}{Proposition}
\title{Complex scratch sharing improves the conditional saving beyond $2^{-33}$}
\author{Research extension prepared with OpenAI Codex}
\date{October 7, 2026}
\begin{document}
\maketitle
\begin{abstract}
We extend the compact-control integer-multiplication draft by sharing
first-stage and third-stage auxiliary roles in its complex network.
The bridge residuals are nondegenerate and nonalternating over $\mathbb F_2$,
so the original translation-kernel phase interface survives. With the same
bit primitive and retained assembly, this yields the conditional bound
$T(n)=O(n(\log n)^{1-\kappa})$ for
$\kappa=125/10^{12}=1.25\cdot10^{-10}>2^{-33}$.
This increases the previous displayed saving $83/10^{12}$ by the factor
$125/83\approx1.506$. The upstream theorem and prior compact-control
extensions are assumed. This new extension has not been independently reviewed.
\end{abstract}

\section{Scope}
The baseline is Douglas Colkitt's \texttt{CrocSwap/integer-mult-bounds},
commit \texttt{6e564879f51ae16f23d392e9e196c605f36d90df}, and its pinned
OpenAI manuscript, commit
\texttt{adc7f1241b42e322a6451854ab7e4b4c146bf78a}.
The model remains a fixed finite alphabet and a fixed finite number of
one-dimensional Turing-machine tapes. We retain the paired bit network at
$h_{\rm b}=50$, the compact-control movement and reservation proofs,
local exceptional repair, independent complex arity, generalized stopped
guard, Gaussian choice, resampling, and final rounding interfaces.
Only the complex scratch allocation and final parameter choices change.
The new argument below is a mathematical derivation, supported by exact
arithmetic and finite checks; it is not a formal or independent verification.

\section{Construction and binary phase proof}
\subsection{Sharing scratch in the complex network}
\label{sec:complex-stage-reuse}
This extension retains every scalar gate and every gate label of the original
complex network at $h=25$. It identifies certain first-stage scratch roles
with third-stage scratch roles. The bit network and all its interfaces remain
unchanged. The following proof supplies the binary residual obligation that
is additional to rational stage sharing.

Write $F=\mathbb F_2^h$, $\mathcal T=\binom{[h]}3$, $v=\binom h3$,
$m=h^3$, and $z=\binom{h-3}3+3(h-3)$. Work with the ordinary dot product
on $F$ and the induced product on $F^{\otimes3}$. Triple indicators have
norm one. Neighbors are distinct triples with even intersection.

\paragraph{A disjoint-image permutation.}
Identify $[h]$ with $\mathbb Z/h\mathbb Z$. Partition triples into translation
orbits. In each orbit choose its lexicographically first representative $T$
and the least positive shift $k$ with $T\cap(T+k)=\varnothing$; apply that
same shift to every member of the orbit. Such a shift exists for $h\ge8$,
because $T-T$ contains at most seven residues. Translation permutes the orbit,
even if it has a stabilizer. The resulting permutation $\pi$ of all triples
therefore satisfies $t_A\cdot t_{\pi(A)}=0$ for every $A$.

\paragraph{Two explicit role matchings.}
Fix a common ordering of triples, and write $n_S(j)$ for the $j$th neighbor
of $S$, $0\le j<z$. A first-stage side role with fixed second and third
triples $(A,B)$ and physical first-coordinate target/source pair
$(S,n_S(j))$ is identified with the third-stage side role having fixed first
and second triples $(S,\pi(A))$ and third-coordinate target/source pair
$(B,n_B(j))$. This is bijective: $A$ is recovered by $\pi^{-1}$, while
$S,B,j$ are read from the third-stage role. Separately, identify central role
$i\in[h]\cup\{*\}$ of the first-stage invocation $(A,B)$ with central
role $i$ of the third-stage invocation $(B,\pi(A))$. This is also bijective.
Second-stage scratch remains separate. No roles used simultaneously are
identified.

Each invocation restores arbitrary scalar scratch, so these identifications
preserve the signed exchange $(X,Y)\mapsto(-Y,X)$ and fix every remaining
scratch role. In particular the argument does not assume zero temporary values.

\paragraph{Nested binary frames.}
The last first-stage label on either type of shared role is
\[
 E=F\otimes\langle t_A\rangle\otimes\langle t_B\rangle,
 \qquad \dim E=h.
\]
For the paired side role the first third-stage label is
\[
 H_{\rm side}=
 \langle t_S\otimes t_{\pi(A)}\rangle^\perp_{F\otimes F}
 \otimes\langle t_B\rangle,
 \qquad \dim H_{\rm side}=h^2-1.
\]
For the paired central role it is
\[
 H_{\rm center}=
 \langle t_B\otimes t_{\pi(A)}\rangle^\perp_{F\otimes F}
 \otimes F,
 \qquad \dim H_{\rm center}=m-h.
\]
The disjointness of $A$ and $\pi(A)$ implies $E\subset H$ in both cases.
All these spaces are nondegenerate: their defining lines have norm one,
and they are tensor products of nondegenerate spaces. Hence the new edge
has nondegenerate orthogonal residual $R=H\cap E^\perp$.

Nondegeneracy alone would not suffice over $\mathbb F_2$. For a side bridge,
choose $x\notin S$ and $y\notin A$. The vector
$e_x\otimes e_y\otimes t_B$ belongs to $H_{\rm side}\cap E^\perp$
and has norm one. For a central bridge choose $x\notin B$, $y\notin A$
and any $q\in[h]$. Then $e_x\otimes e_y\otimes e_q$ belongs to
$H_{\rm center}\cap E^\perp$ and has norm one. Every new residual is
therefore nonalternating and admits an orthonormal basis by the original
binary-space lemma. The side and central residual dimensions are respectively
$h^2-1-h$ and $m-2h$.

All other edges retain their original residuals and orthonormal bases.
For nested labels $H=E\mathbin\perp R$, orthogonal weight additivity gives
\[
 q_H(x)-q_E(x)=\operatorname{wt}(P_Rx)
 =\sum_{u\in\mathcal B_R}\operatorname{wt}(u)[u\cdot x]\pmod4.
\]
Each basis vector has odd weight, so the new edge is implemented by exactly
$\dim R$ forward or inverse translation kernels, just as in the original
complex phase interface. This verifies the phase signs, not only the dimensions.

\paragraph{Exact deficit and endpoints.}
The old two edges on a paired role, $E\to F^{\otimes3}$ and $0\to H$,
have total residual dimension $m-h+\dim H$. Their replacement $E\to H$
has dimension $\dim H-h$, saving exactly $m$ per deleted role.
The total number of deleted roles is
\[
 R_0=v^2(vz+h+1)=19540339540000.
\]
All central decreases are retained, including those in distinct invocations
that now share a role. Thus, writing $N=v^3$ and $L=3v^2h(h+1)$,
\begin{align*}
 W_{\rm c}^{\rm share}&=2N+2v^2(vz+h+1)=39105013080000,\\
 s_{\rm c}^{\rm share}&=W_{\rm c}^{\rm share}m-2N+2L
                       =611015825672000000,\\
 W_{\rm c}^{\rm share}m-s_{\rm c}^{\rm share}&=3703000000,\\
 \eta_{\rm c}^{\rm share}&=\frac1{165005625}.
\end{align*}
Every surviving scratch role still starts with zero label, ends with the
full label, and is fixed by the scalar permutation. Data endpoints and their
weight-27 corrections are unchanged. The original signed-exchange correction
and tensor-column argument consequently give the same completed all-role
phase interface, with the displayed smaller $W,s$.

\begin{proposition}\label{prop:shared-complex-interface}
The shared complex network supports the original phase interface with
$m=15625$, $W=W_{\rm c}^{\rm share}$, $s=s_{\rm c}^{\rm share}$, and
$\sigma=1-627/10^{12}$.
\end{proposition}
\begin{proof}
The frame, residual, and endpoint arguments above prove the interface.
Exact logarithm enclosures certify
$\eta_{\rm c}^{\rm share}>(627/10^{12})\log(15625)$.
Since $e^{-x}>1-x$ for $x>0$, this gives $s/W<m^\sigma$.
\end{proof}

\paragraph{Depth and tape constants.}
No scalar gate is added, and the number of residual kernels decreases.
One can retain the original unshared $h=25$ node-depth bound
$E_{\rm old}=64(W_{\rm old}+m+1)^3$, since its gate schedule is unchanged
and bridges only delete boundary operations. The old $s_{\rm old}$ also
dominates $s_{\rm c}^{\rm share}$. Therefore the entire generalized guard
proof remains valid with $B_{\rm old}=s_{\rm old}+E_{\rm old}$ and its
original $C_0$. The layer recurrence uses the new $W,s$ for time and row
splitting, and these conservative old constants only for coefficient depth.
All matchings, bases, and schedules are fixed finite data, so the retained
fixed-number-of-tapes argument applies without an input-dependent resource.

\section{Full parameter substitution}
Let $a_{\rm b}=296/10^{11}$ and $a_{\rm c}=627/10^{12}$. Choose
\begin{gather*}
 \tau=1-a_{\rm b},\qquad \sigma=1-a_{\rm c},\qquad
 \epsilon=\frac{1999}{10000},\qquad c=\frac14,\\
 \beta=\frac1{1000},\qquad\zeta=\frac1{10000},\qquad
 \delta=\frac1{10^6},\qquad C_1=5-4\beta+\zeta=\frac{49961}{10000},\\
 \lambda=1-\frac{6265}{10^{13}},\qquad
 \lambda'=1-\frac{626}{10^{12}},\qquad
 \kappa=\frac{125}{10^{12}}.
\end{gather*}
The compact-control internal exponent is
$\chi=\tau+(1-\beta)(\sigma-\tau)$, so
\[
 1-\chi=\frac{629333}{10^{15}},\qquad
 1-[\sigma+\beta(1-\sigma)]=\frac{626373}{10^{15}}.
\]
Therefore
\[
 \max\{\tau,\sigma,\chi\}<\lambda<\lambda'<1,\qquad
 \max\{\sigma+\beta(1-\sigma),1-c,0\}<\lambda'.
\]
The guard condition is unchanged: $\epsilon C_1=99872039/10^8<1$.
The reservation exponent $1-c=3/4$ is strictly below $\lambda'$.
Moreover $\epsilon(1+c)=1999/8000<1$, $\epsilon<1/5<1/3$,
$\epsilon c>0$, and $0<\delta<1/8$.

The seven retained assembly margins are exactly
\[
\begin{array}{c|c|c}
 &\text{expression}&\text{value}\\\hline
g_1&1-\epsilon(1+c)&6001/8000\\
g_2&\epsilon c(1-\tau)&73963/(5\cdot10^{14})\\
g_3&\epsilon(1-\lambda')&625687/(5\cdot10^{15})\\
g_4&(1-\tau)(1-\epsilon)&296037/(125\cdot10^{12})\\
g_5&1/4-\delta-5\epsilon/4&31/250000\\
g_6&1-\delta-\epsilon&800099/10^6\\
g_7&\epsilon&1999/10000
\end{array}
\]
Thus
\[
 G=\min_i g_i=\frac{625687}{5\cdot10^{15}}
  =1.251374\cdot10^{-10},\qquad
 G-\kappa=\frac{687}{5\cdot10^{15}}>0.
\]
The strict gap absorbs every retained fixed logarithmic factor.
The old Gaussian and prime comparisons remain unchanged with the same
$\epsilon,\delta$. In particular
$\alpha=\Theta(p^{11999/40000})$ and
$\gamma=O(p^{15997/20000})=o(p)$.
The changed spacing is $K=\Theta(p^{1999/40000})=o(\ell)$, where
$\ell=\Theta(p^{8001/10000})$, and $K/\log p\to\infty$.
The reserved-axis count is $O(\log d+d^{3/4}\log p)=o(d)$.
Record size remains superpolynomial, and all setup and repair costs retain
their original charges. This proves the stated conditional bound.

\section{Reproduction and limits}
The new artifacts are \path{scripts/complex_reuse.py},
\path{certificates/complex-reuse.json},
\path{tests/test_complex_reuse.py}, and
\path{patches/complex-reuse-33.patch}. Run \texttt{make verify}.
The patch applies independently to the unmodified pinned manuscript.
The standalone note is regenerated by
\path{scripts/make_complex_reuse_patch.py}.

Checks enumerate the matching at $h=25$ and other odd and even sizes,
run complete small shared-scratch scalar networks over the rationals and
an odd finite field, build actual binary bridge projectors and orthonormal
bases, verify their phases modulo four, and check all assembly inequalities
using exact fractions. These checks accompany the general proof above;
they do not establish the original multiplication theorem.

Within the same shared-role counting family, exact logarithm enclosures
separate $h=25$ from every integer $22\le h<200$. For $h\ge200$, the
inequalities $\eta<1/(zm)$ and $\log m>1$ give the decreasing upper bound
$1/(zm-1)$ on the recurrence saving, excluding the infinite tail.
This is a limited family comparison, not an optimality claim for other
complex networks or multiplication algorithms. The improvement in $\kappa$
does not imply a comparable measured runtime speedup.
\end{document}
